\chapter{Conclusiones y trabajo futuro}
\label{chap:conclusiones}

\textbf{El primer objetivo queda cumplido.} El sistema aprende una ordenación transversal fuera de
muestra: su Rank-IC medio es 0,1090, supera los baselines deterministas --- el mejor se queda en
0,0130 --- y se mantiene positivo tras placebos, permutación, bootstrap, exclusión de eras, semillas
y neutralización por estilo. En la era reservada, que no intervino en ninguna decisión, el Rank-IC
sigue siendo positivo ($+0{,}0441$). El meta-agente aporta valor frente al promedio equiponderado
porque aprende a reponderar especialistas a partir de cohortes cerradas, si bien el agente de riesgo
por separado ordena algo mejor que la combinación, lo que impide presentar la arquitectura
multi-agente como demostrada por estos datos.

\textbf{El segundo objetivo también queda cumplido, y es el que más margen aportó.} Sobre esa señal
ya congelada, la rejilla de 1.728 configuraciones muestra que las variables de construcción de
cartera afectan al resultado de forma material y muy desigual: el número de posiciones mueve el
Information Ratio mediano en 0,300 y el tope de efectivo en 0,202, mientras que la tolerancia de
deriva y el reparto de pesos son prácticamente inertes (0,016 y 0,009). Optimizar por Information
Ratio eleva el ratio de 0,339 a 0,844 y el exceso geométrico del 2,61 al 6,97\,\%, sin reentrenar
nada. La mejora obtenida sin tocar el modelo es mayor que la acumulada por las tres pasadas
predictivas juntas.

El contraste en la era reservada mide exactamente esa distancia. Con la construcción de cartera
recomendada por el catálogo --- la que estuvo vigente, sin optimizar, durante todo el trabajo
predictivo --- la señal perdía un 11,29\,\% frente al índice, con Information Ratio $-1{,}167$ y
ningún año por encima. Con la cartera elegida por el estudio de optimización --- mismo modelo, misma
señal, mismo panel, sin reentrenar nada --- obtiene un exceso del 2,56\,\% con Information Ratio
$+0{,}304$. El Rank-IC de esa era es positivo en ambos casos, porque no depende de la cartera: la
capacidad de ordenar era idéntica y lo único que cambió fue su gestión.

De ahí se sigue la conclusión que este trabajo defiende: en un sistema de esta clase, la
construcción de cartera no es un detalle de implementación posterior sino una variable capaz de
decidir el signo del resultado fuera de muestra. Ninguna afirmación sobre la utilidad práctica de un
sistema predictivo es interpretable sin declarar con qué cartera se midió, y una evaluación que
reporte únicamente capacidad predictiva puede estar ocultando tanto un éxito como un fracaso.

Nada de esto es una promesa de inversión, y las razones son explícitas: el Deflated Sharpe cae a
0,682 --- encadenar estudios compra Rank-IC al precio de encarecer cualquier afirmación de
rentabilidad ajustada por selección ---, la era reservada aporta seis cohortes cerradas y algo más
de un año de cartera, y la cartera ganadora es la mejor de 1.728 evaluadas. La contribución
defendible del TFM sigue siendo metodológica: demuestra cómo separar señal, selección, robustez y
cartera de forma que preguntas como ésta puedan formularse y responderse sin contaminar la
evidencia.

La prioridad futura no es añadir complejidad al modelo. Es acumular cohortes reservadas suficientes
para someter el hallazgo anterior a una prueba con potencia real, replicar el protocolo en otros
universos y regímenes, y estudiar costes dependientes de liquidez que pongan a prueba una cartera
concentrada en ocho posiciones.

\section{Respuesta a la pregunta de investigación}

La respuesta es afirmativa con condiciones en los dos objetivos.

\textbf{Objetivo 1.} Dentro del universo, periodo y protocolo definidos, los cinco agentes y su meta
producen una ordenación transversal con Rank-IC positivo, superior a los baselines incluidos y
estable frente a contrastes adversariales. El meta parte de pesos iguales y aprende de cohortes
cerradas, por lo que su comportamiento no se reduce a fijar de antemano el especialista ganador.

\textbf{Objetivo 2.} Las variables de construcción de cartera afectan al resultado de forma medible
y desigual, y optimizarlas por Information Ratio produce una cartera claramente mejor que la de
partida. Con esa cartera el sistema supera a SPY en la ventana de selección y también en la era
reservada, neto de los costes modelados; con la cartera de partida no lo hacía en la era reservada.
La respuesta económica es, aun así, mucho más acotada que la predictiva: el Deflated Sharpe y el
tamaño de la muestra efectiva impiden transformar la observación en una promesa de superioridad
ajustada por riesgo, y la cartera elegida es la mejor de 1.728, de modo que su Information Ratio es
una cota superior optimista. La formulación defendible es que el ranking tiene evidencia predictiva
robusta, que su implementación puede ser económicamente compatible con esa señal bajo restricciones
declaradas, y que la distancia entre «puede serlo» y «lo es» la recorre la construcción de cartera y
no el modelo.

La Tabla~\ref{tab:afirmaciones} recoge las cinco preguntas de la introducción convertidas ya en
afirmaciones, cada una con la cifra que la sostiene y con la salvedad sin la cual no debe citarse.
Ninguna de las cinco se sostiene sin su matiz, y esa es precisamente la conclusión de conjunto: hay
evidencia favorable en todas ellas y demostración concluyente en ninguna.

\begin{table}[H]
  \centering
  \footnotesize
  \caption{Las cinco afirmaciones del trabajo, con su evidencia y su salvedad.}
  \label{tab:afirmaciones}
  % Tabla compuesta: las cinco afirmaciones que vertebran el documento, cada una
% con la cifra que la sostiene, el artefacto donde verificarla y su matiz.
\begin{tabularx}{\textwidth}{>{\raggedright\arraybackslash}p{0.24\textwidth}lX}
\toprule
Afirmación & Cifra & Matiz que la acompaña \\
\midrule
El sistema aprende a ponderar a sus especialistas &
$+52\,\%$ &
El meta pasa de reparto uniforme a identificar a \textit{risk} en dos años, pero queda pegado a su
tope de 0,50 desde 2019 \\
\addlinespace
Lo aprendido es señal y no ruido &
$\rankic = 0{,}1004$ &
Frente a 0,0130 del mejor \textit{baseline}, pero 117 cohortes equivalen a $\approx 9$
observaciones independientes \\
\addlinespace
No es suerte &
7 de 8 contrastes &
Permutación $p=0{,}0001$ y placebos en cero; el Deflated Sharpe 0,930 no alcanza su umbral y se
reporta igualmente \\
\addlinespace
Bate al índice en la era reservada &
$+14{,}21\,\%$ &
Exceso geométrico con IR 0,959 y 2 de 2 años, pero el Rank-IC de esa era es indeterminado por falta
de cohortes cerradas \\
\addlinespace
La mejor forma de usar la señal es no interferirla &
$\text{TC}=0{,}247$ &
El perfil que no reordena gana; seis de siete perfiles sesgados destruyen alfa. La cartera
materializa una cuarta parte de la señal \\
\bottomrule
\end{tabularx}

  \caption*{Las cifras predictivas proceden del Model Study de referencia y las económicas de la
  cartera ganadora del Portfolio Study; se desarrollan en los
  Capítulos~\ref{chap:resultados} y~\ref{chap:resultados-economicos}.}
\end{table}

\section{Contribución y agenda}

La contribución principal es un protocolo auditable: separa fecha de publicación y cierre fiscal,
etiqueta futura y entrada de entrenamiento, selección predictiva y cartera, y evidencia de selección
y confirmación reservada. Los artefactos permiten reconstruir configuraciones, pesos, decisiones,
figuras y tablas. Esta disciplina es transferible a otros problemas de ML financiero incluso si el
ranking concreto no se replica.

La segunda contribución es la propia agenda de transferencia, que este trabajo no se limita a
proponer sino que ejecuta: estudiar tamaño de cartera, reparto de pesos, efectivo, tenencia mínima,
suelo de cobertura y rebalanceo sin reabrir la selección predictiva, y hacerlo sobre una serie
recortada para que la ventana reservada siga siendo una comprobación y no un criterio. El resultado
justifica el esfuerzo: la mejora obtenida sin tocar el modelo es mayor que la acumulada por las tres
pasadas predictivas juntas.

El trabajo futuro debe reducir la multiplicidad con un catálogo pre-registrado, esperar etiquetas
reservadas hasta disponer de potencia estadística real, y replicar el protocolo en otros universos y
regímenes. Debe también someter la cartera ganadora a costes dependientes de liquidez, porque una
cartera de ocho posiciones sin efectivo es más sensible a esa especificación que la que se evaluaba
antes. Cada extensión debe preservar la observabilidad temporal que hace interpretables los
resultados actuales.

\section{Agenda ordenada por prioridad}

La evidencia acumulada no sólo señala dónde está el cuello de botella: indica en qué orden conviene
atacar lo que queda.

\begin{enumerate}[leftmargin=*,itemsep=0.4em]
  \item \textbf{Someter la cartera ganadora a una era reservada con potencia.} Es la prioridad
  absoluta, porque de ella depende el hallazgo central. Hoy descansa sobre seis cohortes cerradas y
  algo más de un año de cartera; con tres o cuatro años acumulados podría contrastarse de verdad si
  la optimización generaliza o si su éxito fuera de muestra fue afortunado.
  \item \textbf{Resolver el conflicto entre cadencia y rotación.} La cadencia mensual ganó su
  decisión con una ventaja clara, pero explica algo más de un tercio del flujo de órdenes:
  re-decidir cada mes sobre una señal a doce meses. Separar la frecuencia de \emph{evaluación} de
  la de \emph{ejecución} permitiría conservar las 117 cohortes sin pagar esa parte de la rotación.
  Una fracción comparable procede de la corrección de deriva de pesos, cuya variable apenas mueve
  el Information Ratio, de modo que relajarla es la vía más barata de reducir fricción.
  \item \textbf{Reducir el catálogo antes de volver a ejecutar.} El Deflated Sharpe penaliza por
  todas las configuraciones ensayadas, y encadenar pasadas eleva ese recuento. Un catálogo
  pre-registrado más estrecho, que fije de antemano las variables cuyo valor no se espera que
  cambie, elevaría la probabilidad deflactada sin necesidad de mejorar el modelo.
  \item \textbf{Probar costes dependientes de liquidez sobre la cartera ganadora.} Ocho posiciones
  sin efectivo es una configuración más expuesta al impacto de mercado que la evaluada
  anteriormente, y los costes constantes de esta simulación podrían estar favoreciéndola.
  \item \textbf{Contrastar por qué funciona la microestructura de precio.} La atribución local de la
  Sección~\ref{sec:atribucion-local} descarta la explicación más obvia del predominio de
  \textit{risk} --- no se apoya en la prima de baja volatilidad, sino en huecos y rangos de
  negociación a corto plazo --- pero eso desplaza la pregunta en lugar de cerrarla. Que esas
  variables ordenen retornos a doce meses es un resultado que este trabajo describe y no explica, y
  que merece un diseño específico: podría reflejar una prima de liquidez, un efecto de
  microestructura o un artefacto del universo de gran capitalización.
\end{enumerate}

Obsérvese que ninguna de las cinco prioridades consiste en usar un modelo más potente. La
experiencia de este trabajo apunta en la dirección contraria: la mejora más grande que obtuvo no
vino de mejorar la señal, sino de dejar de desperdiciarla al construir la cartera.

\section{Qué aprendería quien repitiera este trabajo}

Algunas lecciones del desarrollo, documentadas en el Anexo~\ref{anexo:desarrollo}, son
independientes del dominio financiero y probablemente más reutilizables que el resultado empírico.

La primera es que declarar no es computar. Un catálogo puede describir variables que el código nunca
calcula, y el sistema seguirá produciendo resultados plausibles: la comprobación de que lo declarado
coincide con lo computado debe ser automática, porque no es observable a simple vista.

La segunda es que una puerta estadística puede fallar en la dirección contraria a la esperada. La
regla de no inferioridad de este proyecto penalizaba a los candidatos superiores porque su intervalo
era más ancho, y el síntoma no fue un error sino una regularidad sospechosa: nadie resultaba
elegible nunca.

La tercera es que la ausencia de excepción no es evidencia de corrección. Un modelo que devuelve una
constante, una calibración que colapsa, una política de cartera que cobra costes de operaciones no
ejecutadas: todos terminaron su ejecución con éxito aparente y escribieron sus artefactos.

La cuarta es sobre estimadores. Un diagnóstico por deciles basado en la media sugirió durante un
tiempo que la ordenación estaba invertida en todo el histórico; con la mediana, la conclusión se
invertía. Antes de rediseñar un sistema conviene comprobar que el estadístico empleado para
diagnosticarlo sea robusto a la distribución que realmente tienen los datos.

La quinta, y la más incómoda, es que conviene registrar las decisiones tomadas en contra de la
recomendación técnica, con ambos lados del argumento. No mejoran el resultado, pero son
precisamente las que un lector externo no podría reconstruir desde el código.
